\section{Week 5: Uncertainty in Estimates and Fitting Multiple
Participants}

\begin{quote}
\textbf{Chapter 2: Parameter estimation methods} · Lecture 1 h · Lab 2 h
\textbf{Running data:} the Week 2--4 forgetting data (how certain are
our estimates?) and a new simulated \textbf{practice} experiment with 20
participants (the power law of practice).
\end{quote}

\begin{center}\rule{0.5\linewidth}{0.5pt}\end{center}

\subsection{Key terms}

Important terms of this lesson, with their Persian equivalents (repeated
even if they appeared in earlier weeks).

{\def\LTcaptype{none} % do not increment counter
\begin{longtable}[]{@{}
  >{\raggedright\arraybackslash}p{(\linewidth - 4\tabcolsep) * \real{0.3333}}
  >{\raggedright\arraybackslash}p{(\linewidth - 4\tabcolsep) * \real{0.3333}}
  >{\raggedright\arraybackslash}p{(\linewidth - 4\tabcolsep) * \real{0.3333}}@{}}
\toprule\noalign{}
\begin{minipage}[b]{\linewidth}\raggedright
English
\end{minipage} & \begin{minipage}[b]{\linewidth}\raggedright
\RL{\CMfa فارسی}
\end{minipage} & \begin{minipage}[b]{\linewidth}\raggedright
\RL{\CMfa توضیح کوتاه}
\end{minipage} \\
\midrule\noalign{}
\endhead
\bottomrule\noalign{}
\endlastfoot
estimate / estimator & \RL{\CMfa برآورد} / \RL{\CMfa برآوردگر} & \RL{\CMfa مقدار به‌دست‌آمده
برای پارامتر} / \RL{\CMfa روش به دست آوردن آن} \\
fit / fitting & \RL{\CMfa برازش} & \RL{\CMfa تنظیم پارامترها تا مدل با داده جور
شود} \\
parameter & \RL{\CMfa پارامتر} & \RL{\CMfa کمیت قابل تنظیم مدل که بین افراد یا شرایط
فرق می‌کند} \\
simulation & \RL{\CMfa شبیه‌سازی} & \RL{\CMfa تولید داده‌ی مصنوعی با اجرای مدل} \\
power law of practice & \RL{\CMfa قانون توانی تمرین} & \RL{\CMfa کاهش توانی زمان
واکنش با تمرین} \\
exponential / power function & \RL{\CMfa تابع نمایی} / \RL{\CMfa توانی} & \RL{\CMfa دو شکل
رایج برای منحنی فراموشی و یادگیری} \\
sampling distribution & \RL{\CMfa توزیع نمونه‌گیری} & \RL{\CMfa توزیع برآورد در
تکرارهای فرضی آزمایش} \\
standard error (SE) & \RL{\CMfa خطای استاندارد} & \RL{\CMfa پراکندگی برآورد در
تکرارهای فرضی آزمایش} \\
confidence interval (CI) & \RL{\CMfa بازه‌ی اطمینان} & \RL{\CMfa بازه‌ای که با روش
مشخص پارامتر واقعی را در بر می‌گیرد} \\
Hessian & \RL{\CMfa ماتریس} Hessian & \RL{\CMfa ماتریس مشتق‌های دوم؛ انحنای
درست‌نمایی} \\
log-likelihood / NLL & \RL{\CMfa لگاریتم درست‌نمایی} / \RL{\CMfa منفی آن} &
\RL{\CMfa مقداری که بهینه‌سازها کمینه یا بیشینه می‌کنند} \\
Wald interval & \RL{\CMfa بازه‌ی} Wald & \RL{\CMfa برآورد} \ensuremath{\pm} 1.96 SE \\
profile likelihood & \RL{\CMfa درست‌نمایی نیم‌رخ} (profile) & \RL{\CMfa درست‌نمایی یک
پارامتر با بهینه کردن بقیه} \\
likelihood & \RL{\CMfa درست‌نمایی} & \RL{\CMfa احتمال داده به‌عنوان تابعی از
پارامترها} \\
bootstrap & bootstrap \RL{\CMfa (بازنمونه‌گیری)} & \RL{\CMfa تخمین عدم قطعیت با
شبیه‌سازی یا بازنمونه‌گیری} \\
parameter trade-off & \RL{\CMfa بده‌بستان پارامترها} & \RL{\CMfa تغییر یک پارامتر که
با پارامتر دیگر جبران می‌شود} \\
aggregation artifact & \RL{\CMfa خطای ناشی از میانگین‌گیری} & \RL{\CMfa الگویی در
میانگین که در هیچ فردی نیست} \\
trial & \RL{\CMfa کوشش} (trial) & \RL{\CMfa یک بار انجام تکلیف توسط شرکت‌کننده} \\
bounds / constraints & \RL{\CMfa کران‌ها} / \RL{\CMfa قیدها} & \RL{\CMfa محدودیت‌های مجاز
برای مقدار پارامترها} \\
accuracy & \RL{\CMfa دقت} & \RL{\CMfa نسبت پاسخ‌های درست} \\
complete / no / partial pooling & \RL{\CMfa ادغام کامل} / \RL{\CMfa بدون ادغام} /
\RL{\CMfa ادغام جزئی} (pooling) & \RL{\CMfa یک پارامتر برای همه} / \RL{\CMfa جدا برای هر
فرد} / \RL{\CMfa سلسله‌مراتبی} \\
individual differences & \RL{\CMfa تفاوت‌های فردی} & \RL{\CMfa فرق پارامترها بین
افراد} \\
noise & \RL{\CMfa نویز} & \RL{\CMfa تغییرات تصادفی در شواهد یا در داده} \\
benchmark & benchmark \RL{\CMfa (مجموعه‌آزمون استاندارد)} & \RL{\CMfa مجموعه سؤال
برای مقایسه‌ی مدل‌های هوش مصنوعی} \\
\end{longtable}
}

\begin{center}\rule{0.5\linewidth}{0.5pt}\end{center}

\subsection{Learning outcomes}

By the end of this week, students will be able to:

\begin{enumerate}
\def\labelenumi{\arabic{enumi}.}
\tightlist
\item
  \textbf{Explain} the sampling distribution of an estimator and
  \textbf{distinguish} a standard error from a standard deviation of the
  data.
\item
  \textbf{Compute} standard errors and confidence intervals for model
  parameters in three ways: from the Hessian of the NLL (Wald), from the
  profile likelihood, and by parametric bootstrap. \textbf{Compare}
  these three methods.
\item
  \textbf{Diagnose} parameter trade-offs from the estimated correlation
  between parameters.
\item
  \textbf{Demonstrate} by simulation that averaging individual curves
  can produce a functional form (a ``power law'') that no individual
  follows.
\item
  \textbf{Choose} between aggregate, individual and (preview)
  hierarchical fitting, and \textbf{justify} the choice.
\end{enumerate}

\begin{center}\rule{0.5\linewidth}{0.5pt}\end{center}

\subsection{Motivation \& big picture}

\textbf{``\(\hat d = 0.42\)'' is half a result.} Last week the power
model's MLE for the Week 2 data was \(\hat d = 0.42\) (true value 0.5).
Is 0.42 compatible with 0.5? With 0.3? With 0.6? Without a measure of
uncertainty we cannot say whether two participants, two groups or two
conditions really differ. Every parameter estimate we report from now on
comes with a confidence interval.

\textbf{The law that nobody obeys.} The \emph{power law of practice} is
one of the most cited regularities in cognitive psychology: response
times fall with practice approximately as a power function of the number
of trials. Some researchers have argued that the power shape can arise
from \textbf{averaging} over people (or items) who each improve
\emph{exponentially} but at different rates (see Further reading). If
so, the ``law'' describes the average but no single learner, and a
theory built to explain a power function in each individual mind would
be explaining an artifact. This week we demonstrate the effect with
simulated data and learn when to fit individuals rather than group
averages.

\paragraph{Where is this used?}

\begin{itemize}
\tightlist
\item
  \textbf{Any claim of a difference} (``patients forget faster than
  controls'', ``training improved sensitivity'') needs uncertainty
  estimates. Without them a difference of 0.05 in a parameter may be
  pure noise.
\item
  \textbf{Clinical assessment of individuals}: to say that one patient's
  memory decay is abnormal, you need the standard error of \emph{that
  patient's} estimate, not only the group average.
\item
  \textbf{Learning and skill research}: the ``power law of practice''
  was long treated as a universal law. This week's aggregation artifact
  shows why researchers now fit individual learners (and why the same
  caution applies to group-averaged curves in education, sports science
  and human factors).
\item
  \textbf{Replication and meta-analysis} combine estimates from many
  studies by weighting each one by its precision (1/SE²).
\end{itemize}

\begin{center}\rule{0.5\linewidth}{0.5pt}\end{center}

\subsection{Lecture notes}

\paragraph{In plain words}

If you ran the same experiment again on the same person, you would get
slightly different data and therefore a slightly different parameter
estimate. The \textbf{standard error (SE)} says how much the estimate
would jump around between such repetitions; a \textbf{confidence
interval} is a range built from it (roughly estimate \ensuremath{\pm} 2 SE). We cannot
repeat the experiment many times, so we estimate this jumping in three
ways: (1) from how sharply peaked the likelihood is (a sharp peak means
the data pin the parameter down; this is the \textbf{Wald} method using
the \textbf{Hessian}), (2) by tracing how fast the fit gets worse as we
move one parameter (\textbf{profile likelihood}), and (3) by simulating
many fake experiments from the fitted model and refitting each
(\textbf{bootstrap}). The second half of the week warns about a trap:
averaging data across people can create a curve shape that \textbf{no
single person} has. If every learner improves exponentially but at a
different speed, the group average looks like a power law.

\textbf{Suggested timing (60 min):} sampling distributions 10 · Hessian
standard errors 15 · profile likelihood 8 · bootstrap 10 · multiple
participants and aggregation 17.

\subsubsection{1. Estimates have sampling distributions}

If we reran the same experiment on the same participant (same true
\(\boldsymbol\theta^\ast\)), we would get different data and therefore a
different \(\hat{\boldsymbol\theta}\). The distribution of
\(\hat{\boldsymbol\theta}\) over hypothetical replications is its
\textbf{sampling distribution}.

\begin{itemize}
\tightlist
\item
  The \textbf{standard error} (SE) of a parameter is the SD of its
  sampling distribution: how much the \emph{estimate} would vary across
  replications.
\item
  A \textbf{95\% confidence interval} is constructed by a procedure
  that, across replications, contains \(\theta^\ast\) in 95\% of cases.
  (It is \emph{not} ``a 95\% probability that \(\theta^\ast\) is in this
  particular interval''; that is a Bayesian credible interval, Week 9.)
\end{itemize}

We cannot rerun the experiment many times, so we approximate the
sampling distribution in three ways.

\subsubsection{2. Standard errors from the curvature of the likelihood
(Wald)}

\textbf{Intuition.} If the NLL has a sharp, narrow minimum, small
changes in \(\boldsymbol\theta\) make the data much less likely, so the
data pin \(\boldsymbol\theta\) down precisely. A flat, wide minimum
means many values fit almost equally well, so the uncertainty is large.
Curvature is measured by second derivatives.

\textbf{Formalism.} The \textbf{Hessian} of the NLL at the MLE is the
matrix of second derivatives,

\[
\mathbf{H}_{mn} = \frac{\partial^2\, \mathrm{NLL}(\boldsymbol\theta)}{\partial\theta_m\,\partial\theta_n}\bigg|_{\boldsymbol\theta = \hat{\boldsymbol\theta}},
\]

also called the \textbf{observed Fisher information}. For large samples,

\[
\widehat{\mathrm{Cov}}(\hat{\boldsymbol\theta}) \approx \mathbf{H}^{-1}, \qquad
\mathrm{SE}(\hat\theta_m) = \sqrt{[\mathbf{H}^{-1}]_{mm}}, \qquad
\text{95\% Wald CI: } \hat\theta_m \pm 1.96\,\mathrm{SE}(\hat\theta_m).
\]

The off-diagonal elements of \(\mathbf{H}^{-1}\), rescaled to
correlations, show \textbf{parameter trade-offs}: a correlation of −0.9
between two estimates means their errors largely compensate each other
(the tilted valley of Week 3).

\textbf{Sanity check (binomial).} For one proportion,
\(\mathrm{NLL}''(\theta) = y/\theta^2 + (n-y)/(1-\theta)^2\). At
\(\hat\theta = y/n\) this equals \(n / [\hat\theta(1-\hat\theta)]\), so
\(\mathrm{SE} = \sqrt{\hat\theta(1-\hat\theta)/n}\), the familiar
textbook formula.

\textbf{In practice:} compute \(\mathbf{H}\) numerically by finite
differences (lab Part A). The \nolinkurl{hess_inv} returned by BFGS is
only a rough approximation, so do not use it for SEs.

\textbf{When Wald fails:} small samples, estimates near a bound (e.g., a
recall level \(\hat A\) close to 1), strongly non-quadratic likelihoods,
or parameters with skewed sampling distributions (rates, variances).
Then the symmetric \(\pm 1.96\) SE interval can extend into impossible
values.

\begin{quote}
\textbf{Advanced: why this works.} Near its minimum, any smooth NLL is
approximately quadratic (second-order Taylor expansion). A quadratic NLL
is exactly the NLL of a multivariate normal distribution whose inverse
covariance is the Hessian. Combined with the asymptotic normality of the
MLE (Week 4), this gives the result. The same quadratic approximation,
called the \emph{Laplace approximation}, reappears in Bayesian model
comparison (Week 10).
\end{quote}

\paragraph{Reading the math}

\begin{itemize}
\tightlist
\item
  \(\mathbf{H}_{mn} = \partial^2 \mathrm{NLL}/\partial\theta_m\partial\theta_n\)
  at \(\hat{\boldsymbol\theta}\): the matrix of second derivatives,
  i.e.~the \textbf{curvature} of the NLL bowl at its bottom. Diagonal
  entries: how steep the bowl is along each parameter. Off-diagonal
  entries: whether the bowl is tilted (parameters trade off).
\item
  \(\widehat{\mathrm{Cov}}(\hat{\boldsymbol\theta}) \approx \mathbf{H}^{-1}\):
  a sharply curved bowl (large \(\mathbf{H}\)) gives a small covariance,
  i.e.~precise estimates. Taking the inverse is what turns ``steepness''
  into ``uncertainty''.
\item
  \(\mathrm{SE}(\hat\theta_m) = \sqrt{[\mathbf{H}^{-1}]_{mm}}\): take
  the \(m\)-th diagonal entry of the inverse and its square root.
\item
  \textbf{One-parameter check:} for a proportion,
  \(\mathrm{NLL}'' = n/[\hat\theta(1-\hat\theta)]\), so
  \(\mathrm{SE} = \sqrt{\hat\theta(1-\hat\theta)/n}\). With 30 of 50
  recalled: \(\sqrt{0.6 \times 0.4/50} = 0.069\), and the 95\% interval
  is \(0.60 \pm 1.96 \times 0.069 = [0.46, 0.74]\).
\item
  \textbf{Correlation from the covariance:}
  \(\mathrm{corr} = \mathrm{Cov}_{12}/(\mathrm{SE}_1\mathrm{SE}_2)\). A
  value like −0.9 means the two estimates are strongly traded off.
\end{itemize}

\subsubsection{3. Profile likelihood intervals}

For a parameter of interest \(\psi\) (say \(d\)), the \textbf{profile
NLL} is obtained by fixing \(\psi\) on a grid and minimizing the NLL
over all \emph{other} parameters at each grid value:

\[
\mathrm{NLL}_{\text{prof}}(\psi) = \min_{\text{other params}} \mathrm{NLL}(\psi, \text{others}).
\]

The 95\% \textbf{likelihood-ratio (profile) interval} is the set of
\(\psi\) values that are not much worse than the best:

\[
\big\{\psi : 2\,[\mathrm{NLL}_{\text{prof}}(\psi) - \mathrm{NLL}_{\min}] \le 3.84\big\}, \qquad 3.84 = \chi^2_{1,\,0.95}.
\]

Profile intervals follow the actual shape of the likelihood, so they can
be \textbf{asymmetric} and respect natural bounds. They are usually more
accurate than Wald intervals in small samples.

\paragraph{Reading the math}

\begin{itemize}
\tightlist
\item
  \(\mathrm{NLL}_{\text{prof}}(\psi)\): fix the parameter of interest
  \(\psi\) at one value, let the optimizer choose the best values of all
  \emph{other} parameters, and record the resulting NLL. Repeat for many
  values of \(\psi\). The result is a 1-D curve.
\item
  The interval keeps every \(\psi\) whose profile NLL is within
  \(3.84/2 = 1.92\) units of the minimum. The number 3.84 is the 95\%
  point of a \(\chi^2\) distribution with 1 df; it comes from the
  likelihood-ratio test (Week 7).
\item
  Unlike Wald, the curve is not forced to be a symmetric parabola, so
  the interval can be lopsided and never extends to impossible values.
\end{itemize}

\subsubsection{4. The parametric bootstrap}

A simulation answer to ``what if we reran the experiment?'':

\begin{enumerate}
\def\labelenumi{\arabic{enumi}.}
\tightlist
\item
  Fit the model to the data: \(\hat{\boldsymbol\theta}\).
\item
  Simulate \(B\) new datasets from the model at
  \(\hat{\boldsymbol\theta}\), with the same design as the real
  experiment.
\item
  Refit each dataset:
  \(\hat{\boldsymbol\theta}^{*(1)}, \dots, \hat{\boldsymbol\theta}^{*(B)}\).
\item
  The SD of the bootstrap estimates approximates the SE. The 2.5th and
  97.5th percentiles give a \textbf{percentile CI}.
\end{enumerate}

The bootstrap needs no derivatives, works for any statistic (e.g., a
half-life, a difference between conditions), and shows the full shape
and correlations of the sampling distribution. Its cost is \(B\) refits
(\(B \approx 500\)--\(2000\)). The \textbf{non-parametric} bootstrap
resamples the observed trials with replacement instead of simulating
from the model.

\textbf{Comparison:}

{\def\LTcaptype{none} % do not increment counter
\begin{longtable}[]{@{}
  >{\raggedright\arraybackslash}p{(\linewidth - 8\tabcolsep) * \real{0.2000}}
  >{\raggedright\arraybackslash}p{(\linewidth - 8\tabcolsep) * \real{0.2000}}
  >{\raggedright\arraybackslash}p{(\linewidth - 8\tabcolsep) * \real{0.2000}}
  >{\raggedright\arraybackslash}p{(\linewidth - 8\tabcolsep) * \real{0.2000}}
  >{\raggedright\arraybackslash}p{(\linewidth - 8\tabcolsep) * \real{0.2000}}@{}}
\toprule\noalign{}
\begin{minipage}[b]{\linewidth}\raggedright
Method
\end{minipage} & \begin{minipage}[b]{\linewidth}\raggedright
Cost
\end{minipage} & \begin{minipage}[b]{\linewidth}\raggedright
Symmetric?
\end{minipage} & \begin{minipage}[b]{\linewidth}\raggedright
Handles bounds?
\end{minipage} & \begin{minipage}[b]{\linewidth}\raggedright
Needs
\end{minipage} \\
\midrule\noalign{}
\endhead
\bottomrule\noalign{}
\endlastfoot
Wald (Hessian) & One Hessian & Always & No & Smooth NLL, large \(N\) \\
Profile likelihood & One fit per grid value & Can be asymmetric & Yes &
A likelihood \\
Parametric bootstrap & \(B\) fits & Can be asymmetric & Yes & A model
you can simulate \\
\end{longtable}
}

\subsubsection{5. Fitting multiple participants}

Real experiments have \(J\) participants (index \(j\)). There are three
broad strategies:

{\def\LTcaptype{none} % do not increment counter
\begin{longtable}[]{@{}
  >{\raggedright\arraybackslash}p{(\linewidth - 6\tabcolsep) * \real{0.2500}}
  >{\raggedright\arraybackslash}p{(\linewidth - 6\tabcolsep) * \real{0.2500}}
  >{\raggedright\arraybackslash}p{(\linewidth - 6\tabcolsep) * \real{0.2500}}
  >{\raggedright\arraybackslash}p{(\linewidth - 6\tabcolsep) * \real{0.2500}}@{}}
\toprule\noalign{}
\begin{minipage}[b]{\linewidth}\raggedright
Strategy
\end{minipage} & \begin{minipage}[b]{\linewidth}\raggedright
What is fitted
\end{minipage} & \begin{minipage}[b]{\linewidth}\raggedright
Pros
\end{minipage} & \begin{minipage}[b]{\linewidth}\raggedright
Cons
\end{minipage} \\
\midrule\noalign{}
\endhead
\bottomrule\noalign{}
\endlastfoot
\textbf{Aggregate} (complete pooling) & One model to the averaged data &
Stable; works with few trials per person & Assumes everyone is the same;
\textbf{can distort the functional form} \\
\textbf{Individual} (no pooling) & One model per participant; then
analyze the \(J\) sets of estimates & Respects individual differences;
standard practice & Noisy estimates with few trials; second-stage
analysis ignores estimation uncertainty \\
\textbf{Hierarchical} (partial pooling) & Individual parameters drawn
from a group distribution, estimated jointly & Best of both; principled
shrinkage & More complex; usually Bayesian (Week 13) \\
\end{longtable}
}

\textbf{Aggregation artifacts.} The average of nonlinear curves is
generally \textbf{not} a curve of the same family. For example,
averaging exponential curves with different rates does not give an
exponential: fast learners finish early, so the average's relative rate
of change \emph{slows down} over time, which is the defining property of
a power function (Week 2, relative loss rate).

\begin{quote}
\textbf{Advanced: an exact result.} Suppose each participant follows
\(e^{-\lambda_j t}\) and the rates are Gamma-distributed across people,
\(\lambda_j \sim \mathrm{Gamma}(\text{shape } s, \text{rate } b)\). Then
the group average curve is
\[\mathbb{E}\big[e^{-\lambda t}\big] = \Big(1 + \frac{t}{b}\Big)^{-s},\]
which is \emph{exactly} a shifted power function (this is the
moment-generating function of the Gamma distribution, evaluated at
\(-t\)). Everyone forgets exponentially, yet the average forgets by a
power law with exponent equal to the Gamma shape \(s\). You verify this
numerically in Exercise 4.
\end{quote}

\textbf{Practical guideline.} Fit individuals whenever each participant
has enough data to constrain the model; use the aggregate only as a
summary, never as evidence about the \emph{form} of an individual-level
process. When individual data are sparse, use hierarchical models (Week
13).

\paragraph{For an engineer}

\begin{itemize}
\tightlist
\item
  \textbf{Wald SE from the Hessian} = the inverse Fisher information,
  i.e.~the \textbf{CRLB} evaluated at the estimate. MLE reaches the
  bound asymptotically, so for large \(N\) the Wald interval is the
  CRLB-based interval.
\item
  \textbf{The parametric bootstrap} is the Monte Carlo performance
  evaluation you do for an estimator or a receiver (e.g.~simulating many
  noise realizations to get an estimator's variance or a BER curve),
  done at the fitted parameter values.
\item
  \textbf{Profile likelihood} is what you do when you have
  \emph{nuisance parameters}: optimize them out and look at the
  likelihood of the parameter you care about.
\item
  \textbf{Aggregation artifact} = the difference between an
  \textbf{ensemble average} and a \textbf{single realization}. Averaging
  exponential decays with different time constants gives a slower,
  heavy-tailed decay, just as the average of many channel realizations
  does not look like any one realization. The process is not ``ergodic
  across people'': the group mean does not describe the individual.
\item
  \textbf{Partial pooling} (preview of Week 13) is like estimating each
  user's SNR by combining their own few measurements with the
  distribution of SNRs in the whole cell: an empirical-Bayes estimate.
\end{itemize}

\paragraph{For the instructor}

\begin{itemize}
\tightlist
\item
  \textbf{Start with the question ``is 0.42 different from 0.5?''}
  (Motivation). Students will want a yes/no; show that the answer
  depends entirely on the SE.
\item
  \textbf{Show the three methods side by side} on the same data (lab
  Part C table). The message is not ``which is correct'' but ``they
  agree when the likelihood is well-behaved, and differ near bounds or
  with little data''.
\item
  \textbf{The aggregation demo is the highlight.} Plot 5 individual
  exponential curves and their average on a log--log plot: the average
  is almost a straight line (power law), while each individual is
  curved. Ask the class what a researcher who only saw the average would
  conclude.
\item
  \textbf{Typical questions.}

  \begin{itemize}
  \tightlist
  \item
    \emph{``What is the difference between SD and SE?''} SD describes
    the spread of the \emph{data}; SE describes the spread of an
    \emph{estimate} across repeated experiments. SE shrinks with more
    data; SD does not.
  \item
    \emph{``Does a 95\% CI contain the true value with probability
    0.95?''} Not this particular interval; 95\% of intervals built this
    way would. (The Bayesian credible interval of Week 9 does have the
    direct reading.)
  \item
    \emph{``Should I always fit individuals?''} Yes, if each person has
    enough data; otherwise use hierarchical models (Week 13).
  \end{itemize}
\end{itemize}

\begin{center}\rule{0.5\linewidth}{0.5pt}\end{center}

\subsection{Worked example}

\subsubsection{Part 1: three intervals for a recall probability}

A participant recalls \(y = 30\) of \(n = 50\) items;
\(\hat\theta = 0.60\).

\textbf{Wald.} \(\mathrm{SE} = \sqrt{0.6 \times 0.4 / 50} = 0.0693\), so
the 95\% CI is \(0.60 \pm 1.96 \times 0.0693 = [0.464, 0.736]\).

\textbf{Profile (likelihood ratio).} Solve
\(2[\ell(0.6) - \ell(\theta)] = 3.84\) numerically on both sides of 0.6.
The 95\% CI is \([0.462, 0.728]\), slightly \textbf{asymmetric}: it
extends 0.138 below the estimate and 0.128 above, because the likelihood
is skewed toward 0.5, where binomial variance is largest.

\textbf{Interpretation.} With 50 items, a recall probability is only
known to about \ensuremath{\pm}0.13. A difference of 0.10 between two conditions of
this size would not be convincing on its own.

\subsubsection{Part 2: averaging two exponential learners}

Participant 1 forgets at \(\lambda = 0.1\), participant 2 at
\(\lambda = 1.0\). The average curve
\(\bar f(t) = (e^{-0.1t} + e^{-t})/2\):

{\def\LTcaptype{none} % do not increment counter
\begin{longtable}[]{@{}
  >{\raggedright\arraybackslash}p{(\linewidth - 10\tabcolsep) * \real{0.1667}}
  >{\raggedright\arraybackslash}p{(\linewidth - 10\tabcolsep) * \real{0.1667}}
  >{\raggedright\arraybackslash}p{(\linewidth - 10\tabcolsep) * \real{0.1667}}
  >{\raggedright\arraybackslash}p{(\linewidth - 10\tabcolsep) * \real{0.1667}}
  >{\raggedright\arraybackslash}p{(\linewidth - 10\tabcolsep) * \real{0.1667}}
  >{\raggedright\arraybackslash}p{(\linewidth - 10\tabcolsep) * \real{0.1667}}@{}}
\toprule\noalign{}
\begin{minipage}[b]{\linewidth}\raggedright
\(t\)
\end{minipage} & \begin{minipage}[b]{\linewidth}\raggedright
0
\end{minipage} & \begin{minipage}[b]{\linewidth}\raggedright
1
\end{minipage} & \begin{minipage}[b]{\linewidth}\raggedright
2
\end{minipage} & \begin{minipage}[b]{\linewidth}\raggedright
5
\end{minipage} & \begin{minipage}[b]{\linewidth}\raggedright
10
\end{minipage} \\
\midrule\noalign{}
\endhead
\bottomrule\noalign{}
\endlastfoot
\(\bar f(t)\) & 1.000 & 0.636 & 0.477 & 0.307 & 0.184 \\
Relative loss rate \(-\bar f'(t)/\bar f(t)\) & 0.55 & 0.36 & 0.23 & 0.11
& 0.10 \\
\end{longtable}
}

Each participant has a \emph{constant} relative loss rate (0.1 and 1.0),
but the average's rate \textbf{falls} from 0.55 to 0.10, the signature
of a power-like curve. The reason: at later times the average is
dominated by the slow forgetter, whose rate is low.
\textbf{Interpretation:} a researcher fitting only the average would
conclude that ``memories become more resistant with age'', a claim true
of neither participant.

\begin{center}\rule{0.5\linewidth}{0.5pt}\end{center}

\subsection{Lab (2 hours)}

\subsubsection{Goals}

\begin{enumerate}
\def\labelenumi{\arabic{enumi}.}
\tightlist
\item
  Compute Wald, profile and bootstrap intervals for the power model
  fitted to the Week 2 data, and compare them.
\item
  Visualize parameter trade-offs through the estimated correlation and
  the bootstrap cloud.
\item
  Simulate 20 exponential learners, fit exponential and power models to
  the \textbf{average} and to \textbf{each individual}, and interpret
  the contradiction.
\end{enumerate}

\textbf{Timing:} Part A 20 min · Part B 15 min · Part C 25 min · Part D
40 min · Part E 20 min.

\subsubsection{Setup}

\begin{quote}
\textbf{Code · \RL{\CMfa کد:}} \nolinkurl{code/week05.py}, cell 1 · \RL{\CMfa سلول}
1
\end{quote}

\subsubsection{Part A: Wald standard errors from a numerical Hessian (20
min)}

\begin{quote}
\textbf{Code · \RL{\CMfa کد:}} \nolinkurl{code/week05.py}, cell 2 · \RL{\CMfa سلول}
2
\end{quote}

\textbf{Expected output:}

\Needspace{7\baselineskip}
\begin{Verbatim}
   estimate      SE   lower   upper
A    0.8773  0.0626  0.7546  1.0000
d    0.4214  0.0532  0.3171  0.5257
Estimated correlation between A-hat and d-hat: 0.751
\end{Verbatim}

\textbf{Interpret:} both 95\% intervals contain the true values
(\(A^\ast = 0.9\), \(d^\ast = 0.5\)), so the ``miss'' we worried about
in Weeks 3--4 is well within sampling error. Note that the Wald interval
for \(A\) happens to end at 1.000, right at the boundary of possible
values. With slightly noisier data, the symmetric Wald interval would
include impossible values (probabilities above 1), whereas the bootstrap
interval cannot. The positive correlation is the trade-off valley from
Week 3: when the estimate of \(A\) is too high, the estimate of \(d\)
tends to be too high as well.

\subsubsection{\texorpdfstring{Part B: profile likelihood interval for
\(d\) (15 min)}{Part B: profile likelihood interval for d (15 min)}}

\begin{quote}
\textbf{Code · \RL{\CMfa کد:}} \nolinkurl{code/week05.py}, cell 3 · \RL{\CMfa سلول}
3
\end{quote}

\textbf{Expected output:}

\Needspace{4\baselineskip}
\begin{Verbatim}
Profile 95% CI for d: [0.320, 0.525]  (grid step 0.005)
\end{Verbatim}

\subsubsection{Part C: parametric bootstrap (25 min)}

\begin{quote}
\textbf{Code · \RL{\CMfa کد:}} \nolinkurl{code/week05.py}, cell 4 · \RL{\CMfa سلول}
4
\end{quote}

\textbf{Expected output} (takes about 30--60 s):

\Needspace{7\baselineskip}
\begin{Verbatim}
   boot SE    2.5%  97.5%
A   0.0563  0.7724  1.000
d   0.0501  0.3304  0.528
Bootstrap correlation: 0.703
\end{Verbatim}

\textbf{Compare:} put the Wald, profile and bootstrap intervals for
\(d\) side by side. Are they similar? Which is asymmetric? With these
data (8 cells × 50 items) all three methods should agree reasonably
well. They diverge more with fewer items or with estimates near a bound
(Exercise 5).

\subsubsection{Part D: the power law of practice as an aggregation
artifact (40 min)}

\textbf{Design.} \(J = 20\) participants each complete 60 practice
blocks of a visual-search task. Each participant improves
\textbf{exponentially}:

\[
\mathrm{RT}_j(m) = E + B\,e^{-r_j (m-1)} + \varepsilon, \qquad \varepsilon \sim \mathcal{N}(0, 40 \text{ ms}),
\]

with asymptote \(E = 400\) ms, initial gain \(B = 600\) ms, and
individual learning rates
\(r_j \sim \mathrm{Gamma}(\text{shape} = 0.7, \text{scale} = 0.5)\)
(mean 0.35 per block; a skewed distribution, with a few very fast
learners and several very slow ones). Here \(m\) is the practice block,
and \(\varepsilon\) is the noise of a block-mean RT.

\begin{quote}
\textbf{Code · \RL{\CMfa کد:}} \nolinkurl{code/week05.py}, cell 5 · \RL{\CMfa سلول}
5
\end{quote}

\textbf{Expected output:}

\Needspace{8\baselineskip}
\begin{Verbatim}
True learning rates: [1.000e-03 7.000e-03 7.000e-03 2.000e-02 3.000e-02 4.500e-02 1.220e-01
 1.490e-01 1.720e-01 2.360e-01 2.820e-01 3.260e-01 3.870e-01 3.950e-01
 5.340e-01 7.100e-01 7.900e-01 1.002e+00 1.869e+00 2.534e+00]
Average curve, exp: E = 507.5, B = 423.8, r = 0.205, RMSD = 19.26 ms
Average curve, pow: E = 419.2, B = 579.4, r = 0.529, RMSD = 10.16 ms
\end{Verbatim}

\begin{quote}
\textbf{Code · \RL{\CMfa کد:}} \nolinkurl{code/week05.py}, cell 6 · \RL{\CMfa سلول}
6
\end{quote}

\textbf{Expected output (summary line):}

\Needspace{5\baselineskip}
\begin{Verbatim}
Which model fits each participant better? {'exp': 20}
(with SciPy 1.18: {'exp': 19, 'pow': 1}; one participant with a near-zero learning rate is fitted differently)
\end{Verbatim}

\begin{quote}
\textbf{Code · \RL{\CMfa کد:}} \nolinkurl{code/week05.py}, cell 7 · \RL{\CMfa سلول}
7
\end{quote}

\subsubsection{Part E: interpret (20 min)}

\begin{quote}
\textbf{Code · \RL{\CMfa کد:}} \nolinkurl{code/week05.py}, cell 8 · \RL{\CMfa سلول}
8
\end{quote}

\textbf{Expected output:}

\Needspace{8\baselineskip}
\begin{Verbatim}
Mean of true individual rates:      0.481
Mean of estimated individual rates: 0.449 (SD 0.626)
Exponential rate fitted to the average curve: 0.205
Correlation of true and estimated individual rates: 0.971
(with SciPy 1.18: estimated mean 0.457 (SD 0.621), correlation 0.968, because of the same participant)
\end{Verbatim}

\textbf{Write 2--3 sentences for each:} 1. The average curve is better
described by a power function (RMSD about 10 vs.~19 ms), yet all (or,
depending on the SciPy version, 19) of the 20 individuals are better
described by an exponential. Which statement about \emph{learning} is
justified, and which is not? 2. Why is the exponential rate fitted to
the average curve smaller than the mean of the individual rates? (Hint:
which participants dominate the average late in practice?) 3. Individual
fits recover the true rates well here, with 60 blocks per person. What
would change with only 10 blocks per person, and what would you do then?
(Preview of Week 13.)

\subsubsection{Exercises (increasing difficulty)}

\textbf{Exercise 1 \ensuremath{\star}.} Compute Wald 95\% CIs for the \emph{exponential}
model's MLE on the Week 2 data. Do they contain sensible values? Then
compute them at the bogus ``estimate''
\nolinkurl{(A, lam) = (1.0, 0.107)} from last week's
bounded-Nelder--Mead warning. What happens to the Hessian at a point
that is not a minimum?

\textbf{Exercise 2 \ensuremath{\star}.} Rerun the bootstrap with \(B = 50, 200, 1000\)
and different seeds. How stable are the percentile limits? How large
should \(B\) be for a 95\% CI?

\textbf{Exercise 3 \ensuremath{\star}\ensuremath{\star} (a derived quantity).} Use the bootstrap samples
from Part C to compute a 95\% CI for the \textbf{half-life} of the power
model, \(t_{1/2} = 2^{1/d} - 1\). Why is this easy with the bootstrap
and awkward with the Hessian?

\textbf{Exercise 4 \ensuremath{\star}\ensuremath{\star} (the Gamma--power result).} Draw 100 000 rates
\(\lambda \sim \mathrm{Gamma}(\text{shape } s = 0.8, \text{rate } b = 2)\),
average \(e^{-\lambda t}\) over them for \(t \in [0, 30]\), and overlay
the analytic curve \((1 + t/b)^{-s}\) (Lecture §5, Advanced box). Then
fit exponential and power functions to the averaged curve.

\textbf{Exercise 5 \ensuremath{\star}\ensuremath{\star}\ensuremath{\star} (Advanced: coverage).} How often do 95\%
intervals actually contain the true value? Simulate 300 datasets from
the power model (\(A = 0.9\), \(d = 0.5\)) with \(n = 10\) items per
interval, and compute the Wald and profile CIs for \(d\) for each.
Report the coverage (proportion containing 0.5) of each method. Which is
closer to 95\%? Repeat with \(n = 50\).

\subsubsection{Lab check-off}

Show: (1) a table comparing the Wald, profile and bootstrap 95\% CIs for
\(d\); (2) the bootstrap scatter plot with one sentence on the
trade-off; (3) the individual-vs-average comparison and your answer to
interpretation question 1.

\begin{center}\rule{0.5\linewidth}{0.5pt}\end{center}

\subsection{Connection to AI \& LLMs}

\textbf{Benchmark scores have standard errors.} AI models are compared
on benchmarks, and leaderboards often rank models that differ by one or
two percentage points. A benchmark is a \emph{sample} of items, so its
accuracy has sampling error just like a parameter estimate. Two lessons
from this week apply directly: 1. \textbf{Report uncertainty.} A
bootstrap over items gives a confidence interval for the accuracy and
for the difference between two models. 2. \textbf{Respect the design.}
Both models answered the \emph{same} items, so the comparison is paired,
like a within-participant design. Resampling the two models
independently ignores the shared item difficulty and gives intervals
that are too wide.

Item difficulty also varies a lot, which is the item-level version of
the participant heterogeneity in Part D: a single average accuracy hides
it.

\begin{quote}
\textbf{Code · \RL{\CMfa کد:}} \nolinkurl{code/week05.py}, cell 9 · \RL{\CMfa سلول}
9
\end{quote}

\textbf{Expected output:}

\Needspace{6\baselineskip}
\begin{Verbatim}
Accuracy A = 0.74, B = 0.67, difference = 0.07
95% CI, paired bootstrap:   [-0.03, 0.17]
95% CI, unpaired bootstrap: [-0.06, 0.20]
\end{Verbatim}

A 7-point lead on 100 items is not clearly different from zero. A
leaderboard that ranks A above B without an interval is overconfident.

\textbf{Try it (no API needed).} - \ensuremath{\star} How many items would the benchmark
need for the paired interval to exclude 0, if the true difference stays
the same? Find out by simulation. - \ensuremath{\star}\ensuremath{\star} Compute the Wald SE of each
accuracy, \(\sqrt{p(1-p)/n}\), and compare it with the bootstrap SE. -
\ensuremath{\star}\ensuremath{\star} Discussion: a paper reports ``Model A: 74\%, Model B: 67\%'' with no
interval and concludes that A ``reasons better''. List what you would
ask the authors.

\begin{center}\rule{0.5\linewidth}{0.5pt}\end{center}

\subsection{Common mistakes \& misconceptions}

\begin{enumerate}
\def\labelenumi{\arabic{enumi}.}
\tightlist
\item
  \textbf{Confusing SE with SD.} The SD describes the spread of the
  \emph{data}; the SE describes the uncertainty of an \emph{estimate}
  and shrinks as the data grow.
\item
  \textbf{``There is a 95\% probability that the true value is in my
  CI.''} For a frequentist CI, 95\% refers to the long-run behavior of
  the procedure. The probability statement belongs to Bayesian credible
  intervals (Week 9).
\item
  \textbf{Using BFGS's \nolinkurl{hess_inv} for standard errors.} It is
  a crude by-product of the optimizer. Compute the Hessian properly.
\item
  \textbf{Trusting Wald intervals near bounds.} If \(\hat\theta\) is on
  or near a bound, or the interval crosses impossible values, use
  profile or bootstrap intervals.
\item
  \textbf{Reporting parameters without uncertainty.} Comparing two point
  estimates (``participant 1 has \(d = 0.42\), participant 2 has
  \(d = 0.51\)'') is meaningless without intervals.
\item
  \textbf{Inferring individual-level processes from averaged data.}
  Averages of nonlinear curves change functional form (Part D). This
  applies to items, trials and participants alike.
\item
  \textbf{Treating individual estimates as if they were data.}
  Second-stage t-tests on fitted parameters ignore that some estimates
  are much noisier than others. Hierarchical models fix this (Week 13).
\end{enumerate}

\begin{center}\rule{0.5\linewidth}{0.5pt}\end{center}

\subsection{Check your understanding}

\textbf{1.} What is the sampling distribution of an estimator, and why
can't we observe it directly? \textbf{Answer} The distribution of the
estimates we would get across many hypothetical replications of the
experiment with the same true parameters. We have only one dataset, so
we approximate it with the Hessian (asymptotic theory), the profile
likelihood, or simulation (bootstrap).

\textbf{2.} The NLL has a very flat minimum along one direction of the
parameter space. What does that imply for the standard errors and the
correlation of the estimates? \textbf{Answer} Low curvature means a
small Hessian entry in that direction, so a large variance in
\(\mathbf{H}^{-1}\): large SEs. If the flat direction is diagonal, the
estimates are strongly correlated (a trade-off).

\textbf{3.} Derive the Wald SE for a binomial proportion from the second
derivative of the NLL. \textbf{Answer}
\(\mathrm{NLL}''(\theta) = y/\theta^2 + (n-y)/(1-\theta)^2\). At
\(\hat\theta = y/n\) this is
\(n/\hat\theta + n/(1-\hat\theta) = n/[\hat\theta(1-\hat\theta)]\).
Inverting and taking the square root gives
\(\mathrm{SE} = \sqrt{\hat\theta(1-\hat\theta)/n}\).

\textbf{4.} List the steps of a parametric bootstrap for a CI on a model
parameter. \textbf{Answer} Fit the model to get
\(\hat{\boldsymbol\theta}\). Simulate \(B\) datasets from the model at
\(\hat{\boldsymbol\theta}\) using the original design. Refit each one.
Use the SD of the refitted estimates as the SE, and their 2.5th and
97.5th percentiles as the 95\% CI.

\textbf{5.} Why can a profile-likelihood interval be asymmetric while a
Wald interval cannot? \textbf{Answer} The Wald interval assumes a
quadratic NLL (a normal sampling distribution), which is symmetric by
construction. The profile interval uses the actual likelihood, which can
be skewed, e.g., near a bound or for rate parameters.

\textbf{6.} Twenty participants each learn exponentially, with different
rates. The averaged curve is fitted better by a power function. Should
we conclude that learning follows a power law? \textbf{Answer} No.~The
power shape is an aggregation artifact: averaging exponentials with
heterogeneous rates produces a curve whose relative rate of change
decreases over time. Conclusions about the form of individual learning
must come from individual (or hierarchical) fits.

\textbf{7.} Give one advantage and one disadvantage of fitting each
participant separately. \textbf{Answer} Advantage: respects individual
differences and avoids aggregation artifacts. Disadvantage: with few
trials per person, estimates are noisy (and sometimes extreme or on
bounds), and analyzing them in a second stage ignores their different
uncertainties.

\begin{center}\rule{0.5\linewidth}{0.5pt}\end{center}

\subsection{Assignment}

No new graded assignment this week. \textbf{A2 (LS vs.~MLE) is due at
the start of the Week 6 lecture.} Part A of today's lab shows how to add
standard errors to your A2 estimates. Doing so is not required, but it
will strengthen your interpretation.

\textbf{Optional practice (ungraded):} Exercise 5 (coverage) is the best
preparation for Week 7's recovery studies.

\begin{center}\rule{0.5\linewidth}{0.5pt}\end{center}

\subsection{Further reading}

\begin{itemize}
\tightlist
\item
  \textbf{Farrell \& Lewandowsky (2018), chapter 5.} Combining
  information from multiple participants: aggregate vs.~individual
  fitting and the dangers of averaging. \emph{(Main reading.)}
\item
  \textbf{Farrell \& Lewandowsky (2018), chapter 4} (the sections on
  parameter uncertainty and the Hessian), revisited from last week.
\item
  \emph{Optional, outside the course list:} Heathcote, A., Brown, S., \&
  Mewhort, D. J. K. (2000). The power law repealed: The case for an
  exponential law of practice. \emph{Psychonomic Bulletin \& Review, 7},
  185--207. The classic analysis of averaging artifacts in practice
  data.
\end{itemize}

\emph{Preview for Week 6:} Exercise 4 of Week 3 showed a three-parameter
model beating the true model on SSE. Next week we confront this
directly: why good fit is not enough, and how cross-validation measures
what we really want, which is prediction of new data.
